%\chapter{طراحی و معماری حافظهٔ نهان قربانی}
%
%در این فصل به بررسی ساختار و پیاده‌سازی حافظهٔ نهان قربانی می‌پردازیم که راه حلی برای کاهش ترافیک اضافی حافظه است[cite: 1]. این پیمانه\footnote{\lr{Module}} به عنوان یک بافر بسیار کوچک بین حافظهٔ نهان سطح اول و سطوح پایین‌تر قرار می‌گیرد تا از ارسال درخواست‌های پرهزینه جلوگیری کند[cite: 1].
%
%\section{ساختار تمام‌انجمنی و سیاست جایگزینی}
%حافظهٔ نهان قربانی پیاده‌سازی شده در این پروژه دارای یک ساختار تمام‌انجمنی\footnote{\lr{Fully Associative}} است[cite: 1]. بر اساس مشخصات تعریف شده در کدهای شبیه‌ساز، این حافظه به طور پیش‌فرض دارای گنجایش ۸ بلاک\footnote{\lr{Block}} است[cite: 5]. اندازهٔ هر بلاک برابر با ۶۴ بایت در نظر گرفته شده است تا با اندازهٔ بلاک‌های حافظهٔ نهان سطح اول برابری کند و از بروز خطای عدم تطابق در آدرس‌دهی جلوگیری شود[cite: 5]. هر جایگاه در این بافر با استفاده از یک ساختار دادهٔ مشخص مدل‌سازی شده است که شامل برچسب\footnote{\lr{Tag}}، بیت اعتبار\footnote{\lr{Valid Bit}}، داده‌ها و متغیری برای ثبت زمان آخرین دسترسی است[cite: 5].
%
%برای مدیریت سرریز گنجایش محدود این حافظه، از سیاست جایگزینی کمترین استفادهٔ اخیر\footnote{\lr{Least Recently Used (LRU)}} بهره گرفته شده است[cite: 1, 5]. هنگام ورود یک بلاک اخراجی جدید از سطح بالا، سیستم ابتدا بررسی می‌کند که آیا آدرس مورد نظر از پیش در بافر وجود دارد یا خیر[cite: 5]. در صورت نبود آن، به دنبال یک جایگاه خالی می‌گردد و اگر تمام جایگاه‌ها پر باشند، با مقایسهٔ زمان آخرین دسترسی تمام بلاک‌ها، قدیمی‌ترین آن‌ها را انتخاب کرده و بلاک جدید را جایگزین آن می‌کند[cite: 5].
%
%\section{مکانیزم تبادل دوطرفه و مسیریابی پیام‌ها}
%در معماری جدید، منطق مسیریابی پیام‌ها\footnote{\lr{Message Routing}} به گونه‌ای تغییر یافته است که بلاک‌های اخراجی از حافظهٔ نهان سطح اول به جای انتقال به حافظهٔ اصلی، به صورت مستقیم وارد درگاه بالایی حافظهٔ قربانی می‌شوند[cite: 1, 5]. این پیام‌های اخراجی در قالب درخواست‌های نوشتن دریافت می‌شوند و سامانه به محض دریافت آن‌ها، یک پاسخ اتمام نوشتن برای سطح بالاتر ارسال می‌کند تا بافر آن سطح بدون معطلی آزاد شود[cite: 5].
%
%یکی از فرایندهای کلیدی در این معماری، مدیریت جابه‌جایی دوطرفه\footnote{\lr{Swap}} است[cite: 1]. هنگامی که پردازنده به داده‌ای نیاز دارد که در حافظهٔ سطح اول نیست، یک درخواست خواندن به حافظهٔ قربانی ارسال می‌شود[cite: 5]. اگر آدرس درخواستی در این بافر یافت شود (وقوع اصابت)، دادهٔ مورد نیاز به سرعت به سطح اول برمی‌گردد[cite: 1, 5]. این حافظه از سیاست انحصاری\footnote{\lr{Exclusive Policy}} پیروی می‌کند و این بدان معنا است که به محض بازگرداندن داده به سطح بالاتر، نسخهٔ موجود در حافظهٔ قربانی نامعتبر می‌شود تا مالکیت داده به صورت کامل به سطح اول منتقل گردد[cite: 5].
%
%در صورت عدم اصابت، درخواست به صورت ناهمگام\footnote{\lr{Asynchronous}} به درگاه پایینی (حافظهٔ اصلی) فرستاده می‌شود و در یک صف انتظار قرار می‌گیرد[cite: 5]. تمام این فرایندها و تبادل پیام‌ها به کمک ماشین حالت\footnote{\lr{State Machine}} و با بررسی مداوم شرایط ارسال پیام مدیریت می‌شوند تا از بروز چالش بن‌بست\footnote{\lr{Deadlock}} در شبیه‌ساز جلوگیری شود[cite: 1, 5]. افزون بر این، برای مدل‌سازی تاخیر سخت‌افزاری، یک زمان تاخیر دو سیکلی برای پردازش تراکنش‌های ورودی در نظر گرفته شده است که با استفاده از یک شمارنده کاهش می‌یابد تا به صفر برسد و سپس پردازش ادامه یابد[cite: 5].


%\chapter{طراحی و معماری حافظهٔ نهان قربانی}
%میو
%% باید بر اساس فایل victimcache.go این بخش کامل شود
%% سکشن اول مربوط به ویژگی‌های ویکتیم کش، مثل ظرفیت، تاخیر و ... رو به طور کامل توضیح بده. همچنین fsm و توپولوژی و این‌که چطور به هم وصل شدن توضیح داده بشه.
%% در سکشن بعد، همه‌ی موارد، از جمله مکانیزم swap، نحوه‌ی هندل شدن deadlock و به طور کلی توضیح کامل کد مربوط به ویکتیم کش رو بگو.




\chapter{طراحی و معماری حافظهٔ نهان قربانی}

در این فصل به بررسی جزییات طراحی، ساختار سخت‌افزاری و منطق پیاده‌سازی حافظهٔ نهان قربانی\footnote{\lr{Victim Cache}} می‌پردازیم. همان‌طور که پیش‌تر اشاره شد، حافظهٔ نهان قربانی یک بافر کوچک و تمام‌انجمنی\footnote{\lr{Fully Associative}} است که بین حافظهٔ نهان سطح یک و حافظهٔ اصلی قرار می‌گیرد تا از افت عملکرد ناشی از نقصان‌های برخورد جلوگیری کند\footnote{\lr{Conflict Misses}}. این پیمانه در فایل \lr{victimcache.go} پیاده‌سازی شده است که در ادامه ابعاد مختلف آن تشریح می‌شود.

\section{ویژگی‌های سخت‌افزاری، توپولوژی و ماشین حالت حافظهٔ نهان قربانی}

پیکربندی و مشخصات فنی حافظهٔ نهان قربانی در ساختار \lr{VCSpec} و تابع \lr{DefaultVCSpec} تعیین شده است. ویژگی‌های کلیدی این ساختار شامل فرکانس کاری ۱ گیگاهرتز، اندازهٔ بلاک ۶۴ بایت (مطبق با اندازهٔ بلاک در حافظهٔ نهان سطح یک) و تعداد ۸ بلاک فیزیکی برای ذخیره‌سازی داده‌های اخراجی است. تاخیر داخلی این ماژول برابر ۵ چرخهٔ ساعت\footnote{\lr{Latency = 5}} تنظیم شده است که جریمهٔ بررسی بافر قربانی را مدل‌سازی می‌کند.

\subsection{توپولوژی و نحوهٔ اتصال پورت‌ها}

در معماری دارای حافظهٔ نهان قربانی، نحوهٔ اتصال قطعات تغییر می‌کند تا ترافیک خروجی از حافظهٔ نهان سطح یک مستقیما به حافظهٔ اصلی نرود. این توپولوژی در فایل \lr{main.go} به کمک سیستم اتصال مستقیم پیاده‌سازی شده است:
\begin{itemize}
	\item \textbf{پورت بالایی (\lr{TopPort}):} این پورت به عنوان ورودی حافظهٔ قربانی عمل کرده و به پورت پایینی حافظهٔ نهان سطح یک متصل می‌شود. تمامی درخواست‌های خواندن و نوشتنی که از سطح یک خارج می‌شوند (مانند نقصان‌ها و اخراج‌ها) به این پورت وارد می‌شوند.
	\item \textbf{پورت پایینی (\lr{BottomPort}):} این پورت به عنوان خروجی حافظهٔ قربانی به شمار می‌رود و مستقیما به پورت بالای کنترل‌کنندهٔ حافظهٔ اصلی متصل می‌گردد. درخواست‌هایی که در بافر قربانی نیز یافت نشوند (نقصان در قربانی) یا بلاک‌هایی که به دلیل تکمیل ظرفیت باید اخراج شوند، از طریق این پورت به حافظهٔ اصلی ارسال می‌گردند.
\end{itemize}

\subsection{ماشین حالت و مدیریت تراکنش‌ها}

ماژول حافظهٔ نهان قربانی به عنوان یک کامپوننت زمان‌دار\footnote{\lr{Ticking Component}} در هر چرخهٔ ساعت تابع \lr{Tick} را اجرا می‌کند. مدیریت تراکنش‌ها در این ماژول از طریق یک ماشین حالت ناهمگام انجام می‌شود که چرخهٔ حیات هر درخواست را کنترل می‌کند:
\begin{enumerate}
	\item \textbf{دریافت درخواست:} پیام‌های ورودی از طریق \lr{TopPort} دریافت شده و در صف تراکنش‌های بالا\footnote{\lr{vcTopTransaction}} قرار می‌گیرند.
	\item \textbf{مدیریت تاخیر:} تابع \lr{countDown} در هر چرخهٔ ساعت، مقدار شمارندهٔ تاخیر باقی‌مانده\footnote{\lr{CycleLeft}} را برای تراکنش‌های در صف کاهش می‌دهد تا تاخیر سخت‌افزاری ۵ چرخه‌ای اعمال شود.
	\item \textbf{پردازش سر صف:} پس از اتمام تاخیر، تابع \lr{processHeadTransaction} تراکنش را بر اساس نوع آن (خواندن یا نوشتن) به توابع تخصصی هدایت می‌کند تا اصابت یا نقصان بررسی شود.
\end{enumerate}

\section{مکانیزم جابه‌جایی دوطرفه، سیاست جایگزینی و مدیریت ناهمگام}

منطق اصلی حافظهٔ نهان قربانی در فایل \lr{victimcache.go} و از طریق توابع پردازش خواندن، نوشتن و پاسخ‌های حافظه پیاده‌سازی شده است که در ادامه به تفکیک بررسی می‌شوند.

\subsection{ساختار بافر تمام‌انجمنی}

برخلاف حافظهٔ نهان سطح یک که از نگاشت مستقیم استفاده می‌کند، حافظهٔ نهان قربانی به صورت یک آرایه از ساختارهای \lr{VCBlock} پیاده‌سازی شده است که هر مدخل شامل برچسب\footnote{\lr{Tag}}، بیت اعتبار\footnote{\lr{Valid}}، آرایهٔ داده\footnote{\lr{Data}} و زمان آخرین دسترسی است. به دلیل عدم استفاده از تابع نگاشت ریاضی، هنگام جستجوی یک آدرس، کد با استفاده از یک حلقه تمام خانه‌های آرایه را بررسی می‌کند تا اصابت یا نقصان را تشخیص دهد که این دقیقا رفتار سخت‌افزاری یک بافر تمام‌انجمنی را بازسازی می‌کند.

\subsection{مکانیزم جابه‌جایی دوطرفه و خاصیت انحصاری}

هنگامی که پردازنده در حافظهٔ نهان سطح یک دچار نقصان می‌شود اما دادهٔ مورد نظر در حافظهٔ نهان قربانی موجود است (اصابت در قربانی)، فرایند جابه‌جایی دوطرفه\footnote{\lr{Swap}} رخ می‌دهد. این فرآیند در تابع \lr{handleRead} پیاده‌سازی شده است:
دادهٔ درخواستی به سرعت به سمت حافظهٔ نهان سطح یک بازگردانده می‌شود. برای حفظ خاصیت انحصاری\footnote{\lr{Exclusive Property}} و جلوگیری از وجود نسخهٔ تکراری از داده در سطح یک و قربانی، بلافاصله پس از اصابت، بیت اعتبار آن مدخل در قربانی غیرمعتبر می‌گردد (\lr{Valid = false}). این کار باعث می‌شود جایگاه مورد نظر در قربانی آزاد شود تا بلاک اخراجی بعدی بتواند در آن جای گیرد.

\subsection{سیاست جایگزینی کمترین استفادهٔ اخیر}

هنگامی که حافظهٔ نهان سطح یک خطی را اخراج می‌کند و آن را به حافظهٔ قربانی می‌فرستد (\lr{handleWrite})، اگر بافر قربانی کاملا پر باشد، نیاز به یک سیاست جایگزینی است. در این کد از الگوریتم کمترین استفادهٔ اخیر\footnote{\lr{Least Recently Used - LRU}} استفاده شده است:
کد با بررسی فیلد \lr{LastAccessTime} در تمامی بلاک‌های معتبر، قدیمی‌ترین بلاک را شناسایی می‌کند. سپس آن بلاک قدیمی را از طریق \lr{BottomPort} به عنوان یک درخواست نوشتن به حافظهٔ اصلی منتقل می‌کند تا داده‌های تغییریافته از دست نروند. سپس دادهٔ جدید جایگزین آن بلاک می‌شود و زمان دسترسی آن به روز می‌گردد.

\subsection{مدیریت پیام‌های ناهمگام و جلوگیری از بن‌بست}

به دلیل تاخیر بالای حافظهٔ اصلی، سیستم نباید متوقف شود. مدیریت ناهمگام پاسخ‌ها و جلوگیری از بروز بن‌بست\footnote{\lr{Deadlock}} از طریق صف \lr{vcPendingBottom} انجام می‌شود:
وقتی درخواستی در قربانی یافت نمی‌شود (نقصان در قربانی)، یا وقتی یک بلاک به ناچار به سمت حافظهٔ اصلی اخراج می‌شود، درخواست مربوطه همراه با شناسه‌های اصلی آن در صف \lr{pendingBottom} ذخیره می‌گردد و کنترل به شبیه‌ساز بازمی‌گردد. هنگامی که حافظهٔ اصلی پس از گذشت چرخه‌های فراوان پاسخ خود را برمی‌گرداند، توابع \lr{completeBottomRead} و \lr{completeBottomWrite} پیام برگشتی را دریافت کرده، با استفاده از شناسه‌ها مرجع اولیهٔ آن را پیدا می‌کنند و پاسخ را به مقصد درست در بالا هدایت می‌نمایند. در مورد اخراج‌های پس‌زمینه\footnote{\lr{Eviction}} نیز با تنظیم پرچم مربوطه مشخص می‌شود که نیازی به ارسال پاسخ به سمت بالا نیست، که این رویکرد از ایجاد حالت‌های معلق و بن‌بست در شبکه جلوگیری می‌کند.